%% Section 3: DGFT --- Theory and Mechanism. The theoretical core.
%% Theorem statements and proofs are drafted from the protocol invariants
%% actually enforced by the RevTrace runtime (exhaustive online validation:
%% 320 Qwen revision events + 1,048 Whisper commits, 0 violations).

\subsection{Transactions, certificates, and commits}\label{sec:dgft:tx}

DGFT wraps every state projection in a transaction. Throughout this section,
``committed state'' means the committed feedback sub-state $H$ of
Section~\ref{sec:model}; the remaining serving state $\state^{x}$ is never
written by generated output and therefore needs no certification.

\paragraph{Trust model.}
The authority coordinator and versioned ledger are trusted runtime
components; proposers and clients cannot mint or modify certificates. We
assume collision resistance of $h$, atomic certificate consumption with the
ledger compare-and-swap, and immutable session audio bytes once covered by a
certificate. The threat model excludes a Byzantine authority or compromised
ledger. A digest provides content binding, not standalone authentication;
authenticity follows from issuance inside the protected coordinator.

The protocol theorems do not assume a deterministic proposer, greedy
decoding, synchronous GPU kernels, or equality of proposer KV caches. The
authority may use greedy, beam, or sampled decoding provided that the realized
projection and the complete decoding/runtime configuration (including any
randomness policy needed by the contract) are bound by $\chi_j$. Asynchrony is
allowed because the parent-version compare-and-swap serializes installation.
The theorems certify provenance and installation legality, not semantic
correctness of the authority output or repeatability of an independently
rerun stochastic authority.

\begin{definition}[DGFT transaction]\label{def:tx}
A transaction is the object
\[
  \tx_t = \bigl(\mathrm{sid},\, e,\, p,\, v_{\mathrm{parent}},\,
  [a_s,a_e),\, \out_t,\, \proj(\out_t;H_t),\, \mathrm{cert}_t\bigr),
\]
comprising the session id, contract epoch $e$, policy $p$, the parent state
version read, the audio frontier interval covered, the generated output,
its state projection, and the certificate it carries (pending for fine
proposals, issued for authority commits). The proposal granularity
$g_p$ (fine window) and the authority granularity $g_a$ (coarse canonical
segment) define the granularity ratio $\granratio = g_a/g_p \ge 1$.
\end{definition}

\begin{definition}[Certificate]\label{def:cert}
An authority certificate is a binding record
\[
\begin{aligned}
  \cert_j = \mathrm{Bind}\bigl(&\mathrm{sid}_j,e_j,p_j,v_j,
  [f_j,f_{j+1}),\\[-2pt]
  &h(\audio_{[f_j,f_{j+1})}),h(\proj_j^{a}),\chi_j\bigr),
\end{aligned}
\]
where $\proj_j^{a} = \proj(\out_j^{a};H_{v_j})$ is the canonical state
projection computed by the authority against parent version $v_j$, and
$h(\cdot)$ is a content digest over immutable bytes. The configuration
digest $\chi_j$ binds the model checkpoint, tokenizer, decoding parameters,
projection adapter, numerical precision, and authority-runtime version. A
certificate therefore binds
exactly one canonical projection to exactly one commit tuple; it is issued
only when the contract's authority checks pass, and it is single-use:
consumed by the commit that presents it.
\end{definition}

\begin{definition}[Commit validity]\label{def:commit}
An authority commit is a tuple
\[
  \commit_j = (\mathrm{sid}_j,\, e_j,\, p_j,\, v_j,\, [f_j,f_{j+1}),\,
  \proj_j^{a},\, \cert_j).
\]
$\commit_j$ is \emph{valid}, written $\mathrm{Valid}(\commit_j)$, iff
(i) session, contract epoch, and policy match the live session fence;
(ii) $v_j$ equals the currently committed parent version;
(iii) $[f_j,f_{j+1})$ continues the committed frontier without gap or
overlap; and
(iv) certificate binding: $\cert_j$ binds exactly the presented fields, the
digest $h(\proj_j^{a})$ recomputes from the presented projection, and the
certificate has not been consumed before. Only valid commits mutate
committed state.
\end{definition}

Definitions~\ref{def:tx}--\ref{def:commit} realize requirements
R1--R4 (Section~\ref{sec:model:reqs}) structurally: display never produces
a commit (R1), $v_j$ enforced inside the certificate binding gives parent
ownership (R2), $[f_j,f_{j+1})$ enforces frontier monotonicity (R3), and
$\cert_j$ is issued only by authority checks, never by risk signals (R4).
Table~\ref{tab:guarantees} previews how these properties separate DGFT from
weaker protection regimes; the guarantee claims are evaluated in
Section~\ref{sec:eval} (H2, H5).

\tbl{t5_guarantees}

\begin{definition}[Committed provenance ancestry]\label{def:ancestry}
A projection $P$ is an ancestor of committed state $H_j$ iff $P$ labels an
accepted ledger transition on the unique path
$H_0\rightarrow H_1\rightarrow\cdots\rightarrow H_j$. Displayed proposals
and rejected transactions are not ancestors.
\end{definition}

\subsection{Theorem 1: uncertified-ancestry exclusion}\label{sec:dgft:thm1}

\begin{theorem}[Uncertified-Ancestry Exclusion]\label{thm:ancestry}
If committed state is updated exclusively through commits satisfying
$\mathrm{Valid}$ (Definition~\ref{def:commit}), then no fine projection
that is not bound by a valid certificate can become part of the ancestry of
any future committed state in the sense of Definition~\ref{def:ancestry}.
\end{theorem}

\begin{proof}
By induction over the commit sequence $\commit_1,\commit_2,\ldots$.
Base: the initial committed state $H_0$ is contract-defined, so its
ancestry contains no generated output. Step: assume committed state $H_j$
has ancestry containing only certificate-bound projections. The next state
$H_{j+1}$ arises only from a valid $\commit_{j+1}$; by
Definition~\ref{def:commit}(iv), commit application installs exactly the
projection whose digest is bound by $\cert_{j+1}$, since any other content
fails the binding check and is rejected. A fine projection
$\proj(\out^{f};H)$ that differs from the certified canonical projection
$\proj_{j+1}^{a}$ therefore cannot be installed by $\commit_{j+1}$; if a
fine proposal happens to coincide bitwise with $\proj_{j+1}^{a}$, what
enters state is the certified content, attributable to the authority's
certificate rather than to the uncertified fine generation. Provisional
display never mutates committed state by construction of the runtime
(Section~\ref{sec:system}), completing the induction.
\end{proof}

Theorem~\ref{thm:ancestry} is deliberately architecture-agnostic: it depends
on the commit protocol, not on margin accuracy, model determinism,
FP16/FP32 agreement, or the choice of ASR family.

\subsection{Theorem 2: fork-free committed-state linearity}\label{sec:dgft:thm2}

\begin{theorem}[Fork-Free Committed-State Linearity]\label{thm:linear}
Assume commit application is atomic: the validity checks of
Definition~\ref{def:commit} and the state installation execute as one
indivisible compare-and-swap against the committed state, and each
certificate is consumed by at most one commit. Then, under parent-version
checking (Definition~\ref{def:commit}(ii)) and frontier-continuity checking
(Definition~\ref{def:commit}(iii)), the accepted commits of a session form
a fork-free, gap-free, overlap-free linear history of feedback state.
\end{theorem}

\begin{proof}
Atomicity closes the check-then-act race: two distinct commits cannot both
observe the same committed state, both pass the checks, and both install.
Parent-version checking then makes the accepted commits a tree of height at
most one per version: two distinct commits cannot both satisfy
$v_j = \mathrm{version}(H)$ against the same committed state, so the
history is fork-free. Single-use consumption excludes replay of an
already-consumed certificate against a later parent. Frontier-continuity
checking orders accepted intervals $[f_j,f_{j+1})$ into a sequence with
$f_{j+1} = f'_{j}$ for consecutive commits, excluding gaps and overlaps.
Hence the accepted commits are a single linear sequence.
\end{proof}

\begin{remark}[Scope of the concurrency claim]\label{rem:not-linearizable}
Theorem~\ref{thm:linear} proves a serial accepted-commit history with the
ledger compare-and-swap as its serialization point. It is not a claim of
full Herlihy--Wing linearizability for every client invocation and response;
such a claim would require a separate concurrent-object history and
real-time-order specification~\cite{herlihy1990linearizability}.
\end{remark}

\begin{corollary}[Stale-descendant rejection]\label{cor:stale}
A revision computed against a stale parent, i.e.\
$v_{\mathrm{parent}} \neq \mathrm{version}(H_j)$, cannot legally commit.
It fails Definition~\ref{def:commit}(ii) and is rejected without mutating
state.
\end{corollary}

\subsection{Theorem 3: granularity-independent safety}\label{sec:dgft:thm3}

\begin{theorem}[Granularity-Independent Safety]\label{thm:gran}
For any granularity ratio $\granratio \ge 1$, if every authority segment
starts at the committed parent, covers a contiguous audio frontier, yields a
legal state projection, and passes the certificate and version checks of
Definition~\ref{def:commit}, then the conclusions of
Theorems~\ref{thm:ancestry} and~\ref{thm:linear} hold.
\end{theorem}

\begin{proof}
The proofs of Theorems~\ref{thm:ancestry} and~\ref{thm:linear} use only the
certificate binding of Definition~\ref{def:cert}, the validity predicate of
Definition~\ref{def:commit}, and atomic commit application; none of their
steps depends on the size of $[f_j,f_{j+1})$ relative to the fine window.
Varying $\granratio$ changes how much audio a single certificate covers,
not which commits are admissible.
\end{proof}

\begin{corollary}[Same-granularity certification as $\granratio=1$]\label{cor:sg}
Setting $\granratio = 1$ makes proposal and authority granularities
coincide: the directional coverage witness degenerates to the identity and
fine replay degenerates into risk-triggered re-verification. The
same-granularity (SG) configuration is therefore a reduction of DGFT under
the same protocol, not a separate mechanism.
\end{corollary}

\begin{remark}[Canonical-path boundary]\label{rem:canonical}
Corollary~\ref{cor:sg} identifies the protocol, not the arithmetic path: in
our contracts SG runs an FP16 singleton authority while Coarse/Hybrid run an
FP32 batched authority. These are distinct canonical paths; we claim
protocol equivalence, never bitwise identity.
\end{remark}

\subsection{Authority-cost decomposition}\label{sec:dgft:cost}

The safety results above are silent about cost; the following analysis is a
proposition with explicit assumptions, whose predictions we test
(Section~\ref{sec:eval}, H3). Let $B$ be the number of concurrent sessions,
$q$ the SG verification fraction, $r$ the Hybrid selective-replay fraction,
$C_s$ the singleton authority cost per covered region, $C_c(\granratio,B)$
the batched coarse authority cost per segment, $C_{\mathrm{commit}}$ the
transaction/commit overhead, and $Q_s$ the queueing/fragmentation cost of
singleton authority. Then
\begin{align}
  T_{\mathrm{SG}}^{a} &\approx q\,B\,\granratio\,C_s \;+\; Q_s, \label{eq:cost-sg}\\
  T_{\mathrm{DGFT}}^{a} &\approx C_c(\granratio,B) \;+\;
    r\,B\,\granratio\,C_s \;+\; C_{\mathrm{commit}}, \label{eq:cost-dgft}
\end{align}
and end-to-end service time is governed by the slower lane plus repair:
\begin{equation}\label{eq:cost-service}
  T_{\mathrm{service}} \approx \max(T_p, T^{a}) + T_{\mathrm{repair}},
\end{equation}
with proposer and authority running concurrently on separate GPUs.

\begin{proposition}[Model-based operating condition]\label{prop:coarse}
Let end-to-end service time be given by
Eq.~\eqref{eq:cost-service}. Coarse-DGFT improves service throughput over
SG whenever
\[
\begin{split}
  &\max\bigl(T_p,\; C_c(\granratio,B) + rB\granratio C_s +
  C_{\mathrm{commit}}\bigr) + T_{\mathrm{repair}}^{D} \\
  &\quad <\; \max\bigl(T_p,\; qB\granratio C_s + Q_s\bigr)
  + T_{\mathrm{repair}}^{S}.
\end{split}
\]
A simplified sufficient condition, applicable when the authority lane is on
or near the critical path in both configurations ($T^{a} \ge T_p$) and the
repair costs are comparable, is
\[
\begin{split}
  &C_c(\granratio,B) + rB\granratio C_s + C_{\mathrm{commit}} \\
  &\quad <\; qB\granratio C_s + Q_s .
\end{split}
\]
In particular, even a small verification fraction $q$ leaves SG paying
singleton invocation and queueing costs across all $B$ sessions, while a
batched coarse authority amortizes them into $C_c$.
\end{proposition}

Equations~\eqref{eq:cost-sg}--\eqref{eq:cost-dgft} motivate four empirical
hypotheses rather than determine them without fitted costs: (P1) a burst of
singleton authority work can create an SG tail even at 8--10\% verification;
(P2) a batched coarse authority can amortize invocation and queueing costs;
(P3) Hybrid may trade replay fraction $r$ for quality at a performance and
energy cost; and (P4) the ordering can reverse across models because the
complete authority paths---including precision and execution mode---have
model-dependent costs. Section~\ref{sec:eval} tests these hypotheses on the
same workload. Because SG and Coarse/Hybrid use different canonical
arithmetic paths in our implementation (Remark~\ref{rem:canonical}), the
comparison identifies integrated operating points, not a precision-independent
causal effect of $\granratio$ alone.

\subsection{Operating-point trade-offs}\label{sec:dgft:tradeoff}

A deployment selects SG, Coarse, or Hybrid as configuration policy $\pi$,
trading latency, energy, and quality subject to hard safety:
\[
  \max_{\pi}\ J(\pi)
  \quad\text{s.t.}\quad V(\pi)=0,\ C(\pi)\leq B,
\]
where $J$ aggregates the deployment objective, $V$ counts transaction
violations, and $C(\pi)\leq B$ denotes the contract's display and quality
caps.
DGFT supplies the safety-constrained operating envelope; the three evaluated
contracts instantiate three operating points. We do not solve this
optimization online and do not claim any single operating point dominates:
the measured ordering is workload- and model-dependent
(Section~\ref{sec:eval}, H3).
